\chapter{Аналитическая часть}

\section{Задача коммивояжёра}

Задача коммивояжёра~(задача о странствующем торговце)~---~одна из самых
важных задач транспортной логистики, в которой рассматриваются вершины
графа, а также матрица смежности для расстояния между вершинами. Её
суть~---~найти такой порядок посещения вершин графа, при котором пройденное расстояние будет минимально, а каждая вершина будет посещена только один раз.


\section{Алгоритм полного перебора для решения задачи коммивояжёра}

Алгоритм полного перебора для решения задачи коммивояжёра имеет высокую сложность алгоритма ($n!$). Его идея состоит в полном переборе всех возможных путей в графе и выбор наименьшего из них. Такое решение гарантированно даёт самый оптимальный результат. Недостатком алгоритма является его временная сложность, поэтому когда N не является значительно малым, используют эвристические и поисковые алгоритмы~\cite{book_tsp}.

\section{Муравьиный алгоритм}

Муравьиный алгоритм основан на принципе поведения колонии муравьев. Каждый муравей на своём пути оставляет некоторое количество феромона. Количество феромона на дороге влияет на вероятность выбора муравьём этого пути. Таким образом, при большом количестве прохождения муравьёв, наибольшее количество феромона остаётся на оптимальном пути.

Пусть муравей имеет следующие характеристики:
\begin{enumerate}[label=---]
	\item зрение~---~способен определить длину ребра;
	\item память~---~запоминает пройденный маршрут;
	\item обоняние~---~чувствует феромон.
\end{enumerate}

Целевая функция имеет следующий вид:
\begin{equation}
	\label{d_func}
	\eta_{ij} = 1 / D_{ij},
\end{equation}
где $D_{ij}$~---~расстояние из текущего пункта $i$ до заданного пункта $j$.

Формула вычисления вероятности перехода муравья в заданную точку:
\begin{equation}
	\label{possibility_func}
	P_{kij} = \begin{cases}
		\frac{\tau_{ij}^\alpha\eta_{ij}^\beta}{\sum_{q=1}^m \tau^\alpha_{iq}\eta^\beta_{iq}}, \textrm{если вершина ещё не была посещена муравьём k,} \\
		0, \textrm{иначе}
	\end{cases}
\end{equation}
где $\alpha$~---~параметр влияния длины пути~(коэффициент жадности), $\beta$~---~параметр влияния феромона~(коэффициент стадности), $\tau_{ij}$~---~расстояния от города $i$ до города $j$, $\eta_{ij}$~---~количество феромона на пути из города $i$ в город $j$.

Формула зависимости $\alpha$ и $\beta$:
\begin{equation}
	\label{dependency_func}
	\alpha + \beta = 1
\end{equation}

После завершения движения всех муравьёв, феромон испаряется по следующей формуле:
\begin{equation}
\label{update_pheromon_1}
\tau_{ij}(t+1) = (1-p)\tau_{ij}(t) + \Delta \tau_{ij}.
\end{equation}

При этом
\begin{equation}
\label{update_pheromon_2}
\Delta \tau_{ij} = \sum_{k=1}^N \tau^k_{ij},
\end{equation}
где
\begin{equation}
\label{update_pheromon_3}
\Delta\tau^k_{ij} = \begin{cases}
	Q/L_{k}, \textrm{если ребро посещено k-ым муравьем,} \\
	0, \textrm{иначе}
\end{cases}
\end{equation}

Как муравей выбирает путь перечислено ниже.
\begin{enumerate}
\item Каждый муравей имеет список уже посещённых городов (вершин графа), в которые он больше не может ходить.
\item Муравьиное зрение отвечает за желание посетить вершину.
\item Муравьиное обоняние отвечает за ощущение феромона на определённом пути (ребре). При этом количество феромона на пути (ребре) в момент времени $t$ обозначается как $\tau_{i, j} (t)$.
\item После прохождения определённого ребра муравей откладывает на нем некоторое количество феромона, которое показывает оптимальность сделанного выбора (это количество вычисляется по формуле~\eqref{update_pheromon_3}).
\end{enumerate}

